\section{The cubic calculation and one stabilization}\label{sec:cubic-proof}
\subsection{A rational rank-two calculation}

The zero-primary block can be separated without splitting the complementary
eigenlines.  This gives one calculation for three Fano families.

\begin{proposition}[A universal modified residue]
\label{prop:universal-rank-two-residue}
\coverage{fragment}%
\lean{parameterizedRankTwo_finiteReduction,parameterizedRankTwo_normalizedGauge_and_modifiedResidue}%
\uses{lem:A0preserve}%
For \(s=2a+b\), \(sq\ne0\), consider
\[
 z^2\partial_z y=(U_{a,b}+zD)y,\qquad
 U_{a,b}=\begin{pmatrix}
 0&aq&0&a^2q^2\\1&0&bq&0\\0&1&0&aq\\0&0&1&0
 \end{pmatrix},\qquad D=\tfrac12\diag(3,1,-1,-3).
\]
Its characteristic polynomial is \(T^2(T^2-sq)\).  Its zero-primary block
has rank two, nonzero square-zero leading term and canonical modified residue
\begin{equation}
 R_{a,b}=
 \begin{pmatrix}
 -(2a+3b)/(2s)&1\\
 -4a^2/s^2&(b-2a)/(2s)
 \end{pmatrix},\qquad
 \delta^\sharp(a,b)=\frac{4(b-2a)}{b+2a}.
 \label{eq:universal-residue}
\end{equation}
Here \(a,b,q\) are constant under \(\partial_z\).
\end{proposition}

\begin{proof}
The columns of
\[
 C=\begin{pmatrix}
 aq&0&0&-(a+b)q\\
 0&(a+b)q&-aq&0\\
 1&0&0&1\\0&1&1&0
 \end{pmatrix},\qquad \det C=-s^2q^2,
\]
give \(\im U^2\oplus\ker U^2\).  Direct multiplication gives
\[
 J=C^{-1}UC=J_+\oplus N,\qquad
 J_+=\begin{pmatrix}0&sq\\1&0\end{pmatrix},\quad
 N=\begin{pmatrix}0&1\\0&0\end{pmatrix}.
\]
Put \(B=C^{-1}DC\).  The normalized gauge \(G=I+zX+z^2Y+\cdots\),
with zero diagonal blocks in positive orders, starts with
\[
 X_{+0}=\diag(2a/s,2/(qs)),\qquad
 X_{0+}=\diag(-2/(qs),-2a/s).
\]
These matrices solve \(J_iX_{ij}-X_{ij}J_j=-B_{ij}\).
Writing the transformed coefficient as \(J+zB_1+z^2B_2+\cdots\), its
gauge equation gives
\[
 B_1=B+[J,X],\qquad B_2=BX-XB_1-X+[J,Y].
\]
The derivative contributes \(-X\).  Its diagonal blocks vanish, as do
those of \(XB_1\) and \([J,Y]\), so
\[
 (B_1)_{00}=\diag\left(-\frac{2a+3b}{2s},\frac{2a+3b}{2s}\right),
 \qquad ((B_2)_{00})_{21}=(B_{0+}X_{+0})_{21}
 =-\frac{4a^2}{s^2}.
\]
Modification by \(\diag(1,z)\) gives \(R_{a,b}\).  Its characteristic
polynomial is \(T^2+T+(10a-3b)/(4s)\), proving the discriminant formula.
For an original \(z^2F\) term, add \((C^{-1}FC)_{ii}\) to \((B_2)_{ii}\).
Thus only the first off-diagonal gauge coefficient is needed.
\end{proof}

\subsection{The cubic block}
\label{subsec:shared-cubic-packet}

Only a rank-two block can contribute to \(I_{\mathrm{exp}}\).  For a cubic
threefold, the calculation below isolates the unique such candidate and then
computes the residue of its elementary modification.

Let \(X\subset\PP^4\) be a smooth cubic threefold, let \(P=H|_X\), and put
\(q=Q^\ell\) for the line class.  Beauville's formulas
give \cite[main theorem and (2.1)--(2.3)]{Beauville}
\begin{equation}
 P\star P=P^2+6q,\qquad
 P\star P^2=P^3+15qP,\qquad
 P\star P^3=6qP^2+36q^2.
 \label{eq:beauville-cubic-products}
\end{equation}
Thus, in the basis \((1,P,P^2,P^3)\),
\begin{equation}
 U_X=2\begin{pmatrix}
 0&6q&0&36q^2\\1&0&15q&0\\0&1&0&6q\\0&0&1&0
 \end{pmatrix},
 \qquad D_X=\frac12\diag(3,1,-1,-3),
 \label{eq:cubic-small-even-connection}
\end{equation}
and \(z^2\partial_zy=(U_X+zD_X)y\).  This agrees with
\cite[Section~3]{Cai}.

The change of basis \(\diag(1,2,4,8)\) changes \(U_X=2U_{6,15}\)
to \(U_{24,60}\) and leaves \(D_X\) unchanged.  The universal calculation
therefore applies.  Conjugating its residue by \(\diag(2,1)\)
gives the original cubic normalization:
\begin{equation}
 R_X
 =\begin{pmatrix}-19/18&2\\-8/81&1/18\end{pmatrix}.
 \label{eq:RX}
\end{equation}
Its characteristic polynomial is
\((T+1/6)(T+5/6)\).  Thus the two exponent classes differ by
\(2/3\notin\Z\), so the block is counted.  Equivalently,
\begin{equation}
 \boxed{\delta^\sharp(X)=\frac49},
 \qquad I_{\mathrm{exp}}(X)=1.
 \label{eq:cubic-delta}
\end{equation}

\paragraph{From the small connection to the generic block.}
After inverting \(q\), at the small quantum point the rank-two cluster is separated from the
other eigenvalues.  The cyclic-centralizer argument in
Lemma~\ref{lem:cyclic-primary-persistence} continues that cluster over the formal
bulk ring and proves that its centered leading term remains nonzero and
square-zero.  The modified residue is then preserved by flatness.  Thus
the small calculation determines the exponent count at generic quantum
parameters, for every smooth cubic threefold.

\begin{lemma}[Product with one projective line]
\label{lem:special-projective-line}\coverage{absent}%
Let $T$ be a smooth member of the nine families $\mathcal F$ listed in
Theorem~\ref{thm:hodge-conservation}. Suppose at
the small quantum point over the generic Novikov field its even
Euler operator has separated cyclic
primary factors of ranks at most three. On $T\times\PP^1$ each such factor occurs twice with the
same even rank and cyclic centered operator in the full even bulk germ.
A rank-two factor keeps its canonical modified residue discriminant.
Consequently both rank-two counts double.
\end{lemma}
\begin{proof}
The small genus-zero product formula identifies the full connection and
original lattice with the tensor product: Euler multiplication and grading
are sums of those on the factors \cite[formula~(1)]{Behrend}.
The two $\PP^1$ eigenvalues are $\pm2\sqrt{q_f}$. All shifted $T$ clusters
are distinct, since equality between opposite shifts would make the
independent $q_f$ algebraic over the coefficient field of $T$.
The regular splitting of each rank-one $\PP^1$ factor preserves its original
lattice. Tensoring with it adds only scalar connection coefficients, so
rank-two residue discriminants are unchanged.
Lemmas~\ref{lem:cyclic-primary-persistence} and~\ref{lem:triple-primary-persistence} continue each cyclic cluster in
all even bulk coordinates, including the mixed classes $\alpha\otimes p$.
Proposition~\ref{prop:rank2-rigidity} continues the discriminant.
This uses flatness in the transverse directions; a big quantum tensor
formula at mixed bulk is not required.
\uses{prop:generic-spectral-connection-splitting,lem:cyclic-primary-persistence,lem:triple-primary-persistence,prop:rank2-rigidity}%
\imports{Behrend:product}%
\end{proof}

\subsection{The projective endpoint and the contradiction}

Lemma~\ref{lem:special-projective-line} gives
\begin{equation}
 I_{\mathrm{exp}}(X\times\PP^1)=2.
 \label{eq:atomic-product-value}
\end{equation}
On the other hand, over the generic numerical Novikov field
\(\Lambda=\operatorname{Frac}\C[[Q]]\),
\(QH^\bullet(\PP^4)=\Lambda[H]/(H^5-Q)\).  At \(Q\ne0\), multiplication
by \(5H\) has five distinct eigenvalues after algebraic closure.  Thus
\begin{equation}
 I_{\mathrm{exp}}(\PP^4)=0.
 \label{eq:atomic-projective-value}
\end{equation}

\begin{proof}[Proof of Theorem~\ref{thm:every-cubic}]
Proposition~\ref{prop:atomic-lowdim} makes every actual center correction zero
in a weak factorization of smooth projective fourfolds.  Hence
projective weak factorization and
\eqref{eq:numerical-blowup} make \(I_{\mathrm{exp}}\) birationally
invariant.  If \(X\times\PP^1\) were rational, then
\[
 2=I_{\mathrm{exp}}(X\times\PP^1)
 =I_{\mathrm{exp}}(\PP^4)=0,
\]
a contradiction.
\proves{thm:every-cubic}%
\uses{prop:atomic-lowdim,prop:numerical-blowup,lem:special-projective-line,prop:universal-rank-two-residue}%
\imports{AKMW:weak-factorization,Beauville:quantum-products}%
\end{proof}
